\section{Threats to Validity}
\label{sec:threats}

\subsection{Construct Validity}

\paragraph{T1.1 --- ``Promising'' Is Subjective and Unvalidated.}
The qualitative labels (\textit{promising} / \textit{weak\_speculative} /
\textit{drift\_heavy}) in Run~011 were assigned by a single reviewer without
inter-rater reliability testing.
The 15\%$\to$100\% promising rate improvement (Claim~2) is contingent on these labels.
For Subsets~B and~C (Run~013), the promising rate is inferred from filter passage,
not from independent labeling.

\textit{Mitigation}: The filter's generalisation to B and C (0\% drift, 100\% passing)
is independent of the original labels. However, asserting that 39 Subset~B candidates
are all ``promising'' requires independent labeling by domain experts.

\paragraph{T1.2 --- Cross-Domain Status Is Namespace-Based.}
Cross-domain status is determined by the node's \texttt{domain} attribute, not by
semantic domain content.
A node labeled \texttt{chem} representing a metabolite primarily studied in biology
would be miscategorised.

\paragraph{T1.3 --- Rubric Design Dependence.}
The 5-dimension scoring rubric (weights 0.30/0.25/0.20/0.15/0.10) was designed by
the experimenter.
The H4 claim (revised traceability improves deep-candidate ranking) is relative to
the same rubric's naive mode, not to an external ground truth.

\subsection{Internal Validity}

\paragraph{T2.1 --- Semi-Manual Bridge Alignment.}
Bridge entities were domain-expert-selected (NADH, eicosanoids, neurotransmitters),
not discovered algorithmically.
This introduces researcher degrees of freedom: a different experimenter might
select different bridges and obtain different candidate counts.

\textit{Mitigation}: The mechanism (alignment-dependent reachability) is invariant
to bridge choice; the magnitude depends on bridge selection.

\paragraph{T2.2 --- Narrow Scale Range.}
The work tests KGs at 57 nodes (too small for deep cross-domain; Run~008) and
288--536 nodes (Runs~009--013).
Behaviour on production-scale KGs (thousands to millions of nodes) is unknown.

\paragraph{T2.3 --- Curated KGs, Not Production-Quality.}
The Wikidata-derived KGs are manually curated subsets designed to illustrate
known biological pathways.
They differ from production biomedical KGs (e.g., Hetionet, PrimeKG) in scale,
completeness, noise level, and schema diversity.

\subsection{External Validity}

\paragraph{T3.1 --- Three Domain Pairs.}
Reproducibility is tested on 3 independent domain pairs (all bio-chem).
Claim~3 (bridge dispersion predicts yield) is based on 3 data points---insufficient
for statistical characterisation.

\paragraph{T3.2 --- Bio-Chem Domain Pairs Only.}
No non-biological domain pair (e.g., materials science / chemistry) was tested.
Claims 1 and 3 are implicitly scoped to bio-chem domain pairs.

\paragraph{T3.3 --- No Baseline Comparison.}
No existing cross-domain discovery system (SemMedDB, PREDICT, embedding-based)
was evaluated as a baseline.
The paper establishes that the mechanism works, not that it outperforms alternatives.

\subsection{Conclusion Validity}

\paragraph{T4.1 --- No Statistical Testing.}
All comparisons are reported as raw numbers without confidence intervals.
The 15\%$\to$100\% promising rate change is based on 20 labeled candidates.

\paragraph{T4.2 --- Post-Hoc Filter Design.}
The drift filter (Run~012) was designed after inspecting the Run~011 output.
The claimed generalisation is tested in Run~013 but the initial design remains
data-dependent.

\textit{Mitigation}: The filter criteria have a principled justification
independent of the specific candidates: structural/chemical relations describe
molecular facts, not mechanistic hypotheses. This justification was stated
before Run~011 (Run~009 decision memo).

\paragraph{T4.3 --- Top-20 2-Hop Dominance.}
In all subsets and scoring configurations, top-20 candidates are exclusively 2-hop.
Practical users relying on top-20 would not encounter deep cross-domain candidates.
Claim~2's ``useful'' qualifier applies to the filtered set's scientific quality,
not its current top-k visibility.

\subsection{Relation Semantic Insufficiency (Run~014)}

\paragraph{T5.1 --- Class-Level Relations Cannot Validate Substrate-Specific Claims.}
The \texttt{produces} relation between a reaction class and a functional group class
encodes a chemical generalisation that may not hold for the specific substrate in
the path.
Affected candidates: C-1 ($r$\_Methylation $\to$ $fg$\_Piperidine: does not apply
to dopamine methylation in vivo) and B-1/B-2
($r$\_OxidationNat $\to$ $fg$\_Catechol: does not apply to AA as the specific
substrate).
The Run~012 filter cannot detect this failure because \texttt{produces} is
correctly classified as a strong relation in general; the substrate-specific
failure is invisible to a relation-type filter.

\paragraph{T5.2 --- \texttt{is\_substrate\_of} Does Not Encode Causal Transmission.}
The path $A \xrightarrow{\text{is\_substrate\_of}} \mathrm{enzyme}
\xrightarrow{\text{catalyzes}} B$ establishes shared-enzyme membership, not
$A \to B$ causation.
Affected candidate: C-2.

\paragraph{T5.3 --- KG Merges Cell-Type-Specific Knowledge Without Annotation.}
Merging dopaminergic and serotonergic neuron data into a single graph without
cell-type annotation allows compose to generate paths crossing cell-type boundaries.
Such paths may be structurally valid but physiologically inaccessible.
Affected candidate: C-2.

\paragraph{T5.4 --- Path Direction Can Invert Known Pharmacological Relationships.}
The pipeline generates paths in the structural direction allowed by edge orientations.
When the known pharmacological relationship is the reverse of the path direction,
the output is causally misleading.
Affected candidates: B-1 and B-2 (COX $\to$ catechol direction vs.\ the established
catechol $\to$ COX inhibition direction).
